(i) Suppose $f^*$ is closed, $\mathfrak{p}_1\subseteq\mathfrak{p}_2$, and $\mathfrak{q}_1$ contracts to $\mathfrak{p}_1$. The closed set $f^*(V(\mathfrak{q}_1))$ contains $\mathfrak{p}_1$, hence also its specialization $\mathfrak{p}_2$. Thus some $\mathfrak{q}_2\supseteq\mathfrak{q}_1$ contracts to $\mathfrak{p}_2$.

Primes of $B/\mathfrak{q}$ correspond to primes containing $\mathfrak{q}$, and similarly for $A/\mathfrak{p}$. Surjectivity of the quotient spectrum map therefore says exactly that every prime above $\mathfrak{p}$ lifts to one above $\mathfrak{q}$, which is going-up.

(ii) Suppose $f^*$ is open. By \exref{3.26}, the image of $\operatorname{Spec}B_{\mathfrak{q}}$ in $\operatorname{Spec}A$ is $\bigcap_{t\notin\mathfrak{q}}f^*(D(t))$. Each set in this intersection is an open neighborhood of $\mathfrak{p}$, so it contains every prime contained in $\mathfrak{p}$. Hence $\operatorname{Spec}B_{\mathfrak{q}}\to\operatorname{Spec}A_{\mathfrak{p}}$ is surjective.

Primes of $B_{\mathfrak{q}}$ correspond to primes contained in $\mathfrak{q}$, and similarly for $A_{\mathfrak{p}}$. Thus this surjectivity is equivalent to going-down.
